\appendix
\setcounter{section}{0}
\setcounter{theorem}{46}
\section{Appendix: Reductions Between Complete Problems}
\label{sec:equivalences}

In order to prove \autoref{thm:PPAqEquivalentProblems}, we introduce an additional problem that will serve as intermediate problem in our reductions.\\[1mm]


\noindent
\begin{minipage}{\linewidth}
\begin{definition}
	\label{def:leaf-prime} \contour{black}{($\leaf'_q$)}\vspace{-5pt}
	\begin{itemize}[leftmargin=2.15cm,itemsep=4pt]
		\item[\slshape Principle:] Same as $\leaf_q$, but degrees are allowed to be larger (polynomially bounded).
		\item[\slshape Object:] $q$-uniform hypergraph $G = (V, E)$. Designated vertex $v^* \in V$.
		\item[\slshape Inputs:] $\triangleright$ $C:\bit^n \to (\bit^{nq})^k$, with $(\bit^{nq})^k$ interpreted as $k$ many $q$-subsets of $\bit^n$\\
	$\triangleright$ $v^* \in \bit^n$ (usually $0^n$)
		\item[\slshape Encoding:] $V\coloneqq\bit^n$.\\
	For distinct $v_1, \ldots, v_q$, edge $e\coloneqq\set{v_1, \ldots, v_q}\in E$ if $e \in C(v)$ for all $v \in e$
		\item[\slshape Solutions:] $v^*$ if $\deg(v) \equiv 0 \mod{q}$ and\\
	$v \ne v^*$ if $\deg(v) \not\equiv 0 \mod{q}$
	\end{itemize}
\end{definition}
\end{minipage}\vspace{5mm}


\begin{proof}[Proof of \autoref{thm:PPAqEquivalentProblems}]
We show the following inter-reducibilities: (1) $\leaf_q \asymp \leaf_q'$, (2) $\leaf_q' \asymp \bipartite_q$ and (3) $\leaf_q \asymp \lonely_q$.

\medskip

\parasc{(1a) $\boldsymbol{\leaf_q \reducible \leaf_q'}$} Each instance of $\leafq$ is trivially an instance of $\leaf'_q$.

\medskip

\parasc{(1b) $\boldsymbol{\leaf'_q \reducible \leafq}$.} We start with a $\leaf'_q$ instance $(C, v^*)$, where $C$ encode a $q$-uniform hypergraph $G = (V, E)$ with degree at most $k$. Let $t = \ceil{k/q}$. We construct a $\leaf_q$ instance encoding a hypergraph $\overline{G} = (\overline{V}, \overline{E})$ on vertex set $\overline{V}\coloneqq V \times [t]$, intuitively making $t$ copies of each vertex.

In order to locally compute hyperedges, we first fix a canonical algorithm that for any vertex $v$ and any edge $e \in E$ incident on $v$, assigns it a label $\ell_v(e) \in [t]$, with at most $q$ edges mapping to the same label --- e.g. sort all edges incident on $v$ in lexicographic order and bucket them sequentially in at most $t$ groups of at most $q$ each. Note that we can ensure that for any vertex $v$ at most one label gets mapped to by a non-zero, non-$q$ number of edges. Moreover, if $\deg(v) \equiv 0 \mod{q}$, then exactly $q$ or $0$ edges are assigned to any label.

We'll assume that $\deg(v^*) \not\equiv 0 \mod{q}$, as otherwise, a reduction wouldn't be necessary. We let $(v^*, \ell^*)$ be the designated vertex of the $\leaf_q$ instance, where $\ell^*$ is the unique label that gets mapped to by a non-zero, non-$q$ number of edges incident on $v^*$.

For any vertex $(v, i) \in \overline{V}$, we assign it at most $q$ edges as follows: For each edge $e = \set{v_1, \ldots, v_q}$ such that $\ell_v(e)=i$, the corresponding hyperedge of $(v,i)$ is ${(v_1, \ell_{v_1}(e)), \ldots, (v_q, \ell_{v_q}(e))}$. It is easy to see that the designated vertex $(v^*, \ell^*)$ indeed has non-zero, non-$q$ degree. Moreover, a vertex $\deg(v,i) \notin \set{0, q}$ in $\overline{G}$ only if $v$ has a non-multiple-of-$q$ degree in $G$. Thus, solutions to the $\leaf_q$ instance naturally maps to solutions to the original $\leaf_q'$ instance.

By \autoref{rem:simplifications}, this completes the reduction since the edges are locally computable with black-box access to $C$ and $\overline{V}$ is efficiently indexable.

\medskip

\parasc{(2a) $\boldsymbol{\leaf'_q \reducible \bipartiteq}$.}
We start with a $\leaf'_q$ instance $(C, v^*)$, where $C$ encode a $q$-uniform hypergraph $G = (V, E)$. We construct a $\bipartite_q$ instance encoding a graph $\overline{G} = (\overline{V} \cup \overline{U}, \overline{E})$ such that $\overline{V} = V$ and $\overline{U} = \binom{V}{q}$, i.e. all $q$-sized subsets of $V$. We include the edge $(v,e) \in \overline{E}$ if $e \in E$ is incident on $v$. The designated vertex for the $\bipartite_q$ instance is $v^*$ in $\overline{V}$.

Clearly, all vertices $e \in \overline{U}$ have degree either $q$ or $0$. For any $v \in \overline{V}$, the degree of $v$ in $\overline{G}$ is same as its degree in $G$. Thus, any solution to the $\bipartite_q$ instance immediately gives a solution to the original $\leaf'_q$ instance. By \autoref{rem:simplifications}, this completes the reduction since the edges are locally computable with black-box access to $C$ and $\overline{V}$ and $\overline{U}$ are efficiently indexable (cf. \cite[\S2.3]{r31} for efficiently indexing $\overline{U}$).

\medskip

\parasc{(2b) $\boldsymbol{\bipartiteq \reducible \leaf'_q}$.} We start with a $\bipartite_q$ instance $(C, v^*)$ encoding a bipartite graph $\calG = (V \cup U,E)$ with maximum degree of any vertex being at most $k$. We construct a $\leaf'_q$ instance encoding a hypergraph $\overline{G} = (\overline{V}, \overline{E})$ such that $\overline{V} = V$ with designated vertex $v^*$.

First, we fix a canonical algorithm that for any vertex $u \in U$ with $\deg_{G}(u) \equiv 0 \mod{q}$ produces a partition of it's neighbors with $q$ vertices of $V$ in each part. Now, the set of $q$-uniform hyperedges incident on any vertex $v \in \overline{V}$ in $\overline{E}$ can be obtained as: for all neighbors $u$ of $v$, with $\deg_G(u) \equiv 0 \mod{q}$, we include a hyperedge consisting of all vertices in the same partition as $v$ among the neighbors of $u$ (we ignore neighbors $u$ with $\deg(u) \not\equiv 0 \mod{q}$).

Observe that $\deg_{\overline{G}}(v) \le \deg_{G}(v)$ and equality holds if and only if all neighbors of $v$ in $G$ have degree $\equiv 0 \mod{q}$. Hence for any $v \in \overline{V}$, if $\deg_{\overline{G}}(v) \ne \deg_G(v) \mod{q}$, then there exists a neighbor $u \in U$ of $v$ in $G$ such that $\deg(u) \not\equiv 0 \mod{q}$. Thus, if $v = v^*$ and $\deg_{\overline{G}}(v^*) \equiv 0 \mod{q}$, then either $\deg_G(v) \equiv 0 \mod{q}$ or we can find a neighbor $u$ of $v$ in $G$ with $\deg(u) \not\equiv 0 \mod{q}$. Similarly if for some $v \ne v^*$, we have $\deg_{\overline{G}}(v^*) \not\equiv 0 \mod{q}$, then either $\deg_G(v) \not\equiv 0 \mod{q}$ or we can find a neighbor $u$ of $v$ in $G$ with $\deg(u) \not\equiv 0 \mod{q}$. Thus, any solution to the $\leaf'_q$ instance gives us a solution to the original $\bipartite_q$ instance. This completes the reduction since $\overline{V} = \bit^n$ and the edges are locally computable with black-box access to $C$.

\medskip

\parasc{(3a) $\boldsymbol{\leaf_q \reducible \lonely_q}$.} We start with a $\leaf_q$ instance $(C, v^*)$, where $C$ encode a $q$-uniform hypergraph $G = (V, E)$ with degree at most $q$. If $\deg_G(v^*) = q$ or $0$, then we don't need any further reduction. Else, we construct a $\lonely_q$ instance encoding a $q$-dimensional matching $\overline{G} = (\overline{V}, \overline{E})$ on vertex set $\overline{V} = V \times [q]$. The designated vertices will be $V^* = \set{(v,q-i) : 1 \le i \le q - \deg(v^*)}$. Note that, $|V^*| = q - \deg_G(v^*)$ and hence $1 \le |V^*| \le q-1$.

In order to locally compute hyperedges, we first fix a canonical algorithm that for any vertex $v$ and any edge $e \in E$ incident on $v$, assigns it a unique label $\ell_v(e) \in [q]$ --- e.g. sort all edges incident on $v$ in lexicographic order and label them sequentially in $[q]$. In fact, we can ensure that an edge incident on $v$ get labeled within $\set{1, \ldots, \deg_G(v)}$.

For any vertex $(v, i) \in \overline{V}$, we assign it at most one hyperedge as follows:
\begin{itemize}[label=$\triangleright$]
\item If $\deg_G(v) = 0$, we include the hyperedge $\set{(v,i) : i \in [q]}$.
\item Else if $\deg_G(v) \ge i$, then for edge $e = \set{v_1, \ldots, v_q}$ incident on $v$ such that $\ell_v(e)=i$, the corresponding hyperedge of $(v,i)$ is ${(v_1, \ell_{v_1}(e)), \ldots, (v_q, \ell_{v_q}(e))}$.
\item Else if $0 < \deg_G(v) < i$, we leave it isolated.
\end{itemize}

It is easy to see that our definition of hyperedges is consistent and that the designated vertices $V^*$ are indeed isolated. Moreover, a vertex $(v,i)$ is isolated in $\overline{G}$ only if $1 \le \deg_G(v)\le q-1$. Thus, solutions to the $\leaf_q$ instance naturally maps to solutions to the original $\leaf_q'$ instance.

By \autoref{rem:simplifications}, this completes the reduction since the edges are locally computable with black-box access to $C$ and $\overline{V}$ is efficiently indexable.

\medskip

\parasc{(3b) $\boldsymbol{\lonely_q \reducible \leaf_q}$.} We start with a $\lonely_q$ instance $(C, V^*)$, where $C$ encode a $q$-dimensional matching $G = (V, E)$. We construct a $\leaf_q$ instance encoding a $q$-uniform hypergraph $\overline{G} = (\overline{V}, \overline{E})$ on vertex set $\overline{V}$ that will be specified shortly. We describe the hyperedges in $\overline{G}$ and it'll be clear how to compute the hyperedges for any vertex locally with just black-box access to $C$.

We start with $\overline{V} = V$. Our goal is to transform all vertices of degree $1$ to degree $q$, while ensuring that vertices of degree $0$ are mapped to vertices of degree not a multiple of $q$. Towards this goal we let $\overline{E}$ to be set of edges in $E$ in addition to $q-1$ canonical $q$-dimensional matchings over $V$. For example, for a vertex $v \coloneqq (x_1, \ldots, x_n) \in V = [q]^n$, the corresponding edges in $\overline{E}$ include an edge in $E$ (if any) and edges of the type $e_i = \set{(x_1, \dots, x_{i-1}, j, x_{i+1}, \dots, x_n) : j \in [q]}$ for $i \in [q-1]$ (note, this requires us to assume $n \geq q-1$). Adding the $q-1$ matchings increases the degree of each vertex by $q-1$. Therefore, vertices with initial degree $1$ now have degree $q$ and vertices with initial degree $0$ now have degree $q-1$. However, a couple of issues remain in order to complete the reduction, which we handle next.

{\em Multiplicities.} An edge $e \in E$ might have gotten added twice, if it belonged to one of the canonical matchings. To avoid this issue altogether, instead of adding edges directly on $V$, we augment $\overline{V}$ to become $\overline{V} \coloneqq V \cup \inparen{\binom{V}{q} \times [q-1]}$, i.e. in addition to $V$, we have $q-1$ vertices for every potential hyperedge of $G$. For any edge $e\coloneqq\set{v_1, \ldots, v_q} \in E$, instead of adding it directly in $\overline{G}$, we add hyperedge $\set{v, (e,1), (e,2), \ldots, (e,q-1)}$ for each $v \in e$. Note that, all vertices $(e,i) \in \binom{V}{q} \times [q-1]$ have degree $q$ if $e \in E$ and degree $0$ if $e \notin E$, so they are non-solutions for the $\leaf_q$ instance. For vertices in $V$, we still have as before that vertices with initial degree $1$ now have degree $q$ and vertices with initial degree $0$ now have degree $q-1$.

{\em Designated vertex.} In a $\leaf_q$ instance, we need to specify a single designated vertex $v^* \in \overline{V}$. If the $\lonely_q$ instance had a single designated vertex then we would be done. However, in general it is not possible to assume this (for non-prime $q$). Nevertheless, we provide a way to get around this. We augment $\overline{V}$ with $t = (q-1)(q-k) + 1$ additional vertices to become $\overline{V} \coloneqq V \cup  \inparen{\binom{V}{q} \times [q-1]} \cup \set{w_{i,j} : i \in [q-k],\ j \in [q-1]} \cup \set{v^*}$, where $v^*$ will eventually be the single designated vertex for the $\leaf_q$ instance.

Let $V^* = \set{u_1, \ldots, u_k} \subseteq V$ be the set of designated vertices in the $\lonely_q$ instance (note $1 \le k < q$). So far, note that $\deg_{\overline{G}}(u_i) = q-1$. The only new hyperedges we add will be among $u_i$'s, $w_{i,j}$'s and $v^*$, in such a way that $\deg_{\overline{G}}(u_i)$ will become $q$, the degree of all $w_{i,j}$'s will also be $q$ and degree of $v^*$ will be $q-k$.
\begin{itemize}[label=$\triangleright$, itemsep=3pt]
\item For each $u \in V^*$, include $\set{u, w_{1,1}, \ldots, w_{1,q-1}}$. So far, $\deg_{\overline{G}}(u) = q$ and $\deg_{\overline{G}}(w_{1,j}) = k$.
\item For each $j \in [q-1]$ and each $i \in \set{2, \ldots, q-k}$, include $\set{w_{1,j}, w_{i,1}, \ldots, w_{i,q-1}}$.\\
So far, $\deg_{\overline{G}}(w_{i,j}) = q-1$ for all $(i,j) \in [q-k] \times [q-1]$.
\item Finally, for each $(i,j) \in [q-k] \times [q-1]$, include $\set{v^*, w_{i,1}, \ldots, w_{i,q-1}}$.\\
Now, $\deg_{\overline{G}}(w_{i,j}) = q$ for all $(i,j) \in [q-k] \times [q-1]$ and $\deg_{\overline{G}}(v^*) = q-k$.
\end{itemize}

Thus, we have finally reduced to a $\leaf_q$ instance encoding the graph $\overline{G} = (\overline{V}, \overline{E})$ with $\overline{V} \coloneqq V \cup  \inparen{\binom{V}{q} \times [q-1]} \cup \set{w_{i,j} : i \in [q-k],\ j \in [q-1]} \cup \set{v^*}$. By \autoref{rem:simplifications}, this completes the reduction, since $\overline{V}$ is efficiently indexable (again, see \cite{r31} for a reference on indexing $\binom{V}{q}$) and the edges are locally computable using black-box access to $C$.
\end{proof}

